\section{Conclusion}
\label{sec:conclusion}

This paper presented a hybrid architecture for prioritized multi-AGV
scheduling in a multi-floor smart-manufacturing layout. The work began with
an implementation of the prioritized-planning baseline of Guo et al.
(ICARCV 2024) \cite{guo2024}, correcting the original demonstration's documented issues in
conflict resolution, parking-node handling, service-time placement, and
parameterization. The corrected baseline assigns jobs in earliest-pickup
order using a fitness score, constructs shortest-path routes, and schedules
AGVs by fixed priority while eliminating conflicts through sliding-window
delays.

The main architectural contribution divides planning between a central
server and one server per floor. The central server assigns jobs, separates
routes into floor legs, and books shared lift trips. Floor servers construct
conflict-free schedules for their own legs using local node and track
reservations. Lift-lobby waiting points let an AGV wait near the floor-level
conflict or lift that is delaying it. An independent checker validates each
complete plan for route endpoints, node and bidirectional-track conflicts,
edge travel times, dwell windows, and exactly-once job assignment. On the
tested seeds, the hybrid planner in priority order reproduces the centralized
planner with lobbies, providing a check that the route decomposition
preserves that reference schedule under that processing order.

In the saved comparison of 100 jobs and 20 AGVs over ten seeds, the default
FCFS hybrid configuration reduced mean last-delivery makespan by 65\%, total
tardiness by 72\%, and load imbalance by 80\% relative to the paper-style
centralized configuration. The busiest-server CPU time fell by 92\%, and the
maximum number of reservation windows held by one server fell by 69\%. These
results demonstrate the potential value of lobby waiting and floor-level
decomposition for the modeled three-floor factory. They are empirical
results for the tested layout and settings, not a general guarantee of
optimality or of the same improvement at larger scales \cite{projectRepo}.

The reduction in tardiness is substantial, but the mean total tardiness for
the hybrid remains approximately 90,105~s across 100 jobs per seed. This
indicates that conflict resolution and decomposition alone do not ensure
that assigned deliveries meet their desired times. The current job assignment
is greedy and priority-dependent, and delay accumulated by shared routes,
lift use, and safety separation can still move deliveries well past their
targets. Reducing this residual tardiness will require changes to assignment,
sequencing, and resource coordination in addition to distributing the
planning workload.

\paragraph{Future work.}
Future work will extend the static, known-job setting to dynamic job arrivals
and investigate battery limits and charging decisions. The lift model can
be generalized to multiple lifts per floor, with lift capacity and service
policies represented explicitly. A further architectural step is to remove
the central lift-booking server and investigate fully decentralized lift
coordination. Broader comparisons should include auction-based assignment,
deep reinforcement learning (DRL), and conflict-based search (CBS), using
common instances, objectives, and verification rules. The design should also
be validated on physical multi-floor hardware.

